% FAC-EN-U0067
% Source: fac_body.tex lines 2524--2587.
% Authority: PDF physical pp. 67--68; printed pp. 262--263
% (Chapter III no. 69 complete).
% French control: T0040 repairs eight inherited epsilon commands to the
% printed membership sign. No printed-source correction is admitted here.

\phantomsection\label{fac:III:4:no:69}
\textbf{69. Interpretation of $H_k^q(M)$ by means of
$\operatorname{Ext}_S^q$.}
Let $M$ be a graded $S$-module, and let $k$ be an integer $\geq0$. Put
\phantomsection\label{fac:III:4:no:69:def:Bk}
\[
B_k^q(M)=\sum_{n\in\mathbf{Z}}H_k^q(M(n)),
\]
with the notation of \hyperref[fac:III:3:no:61]{no.~61}.

This gives a graded group, isomorphic to the $q$-th cohomology group of
the complex $\sum_{n\in\mathbf{Z}}C_k(M(n))$. This complex may be given
an $S$-module structure compatible with its grading by putting
\phantomsection\label{fac:III:4:no:69:formula:S-action}
\[
(P\cdot m)\langle i_0\cdots i_q\rangle
=P\cdot m\langle i_0\cdots i_q\rangle,
\quad\text{if }P\in S_p\text{ and }
m\langle i_0\cdots i_q\rangle\in C_k^q(M(n)).
\]
Since the coboundary operator is a homogeneous $S$-homomorphism of
degree~$0$, it follows that the $B_k^q(M)$ are themselves graded
$S$-modules.

We shall put
\phantomsection\label{fac:III:4:no:69:def:Bq}
\[
B^q(M)=\varinjlim_{k\to\infty}B_k^q(M)
=\sum_{n\in\mathbf{Z}}H^q(M(n)).
\]
The $B^q(M)$ are graded $S$-modules. For $q=0$ we have
\phantomsection\label{fac:III:4:no:69:formula:B0}
\[
B^0(M)=\sum_{n\in\mathbf{Z}}H^0(M(n)),
\]
and recover the module denoted by
\hyperref[fac:III:3:no:67:def:M-natural]{$M^\natural$} in
\hyperref[fac:III:3:no:67]{no.~67} (when $M$
\hyperref[fac:III:2:no:56:def:TF]{satisfies condition~(TF)}).
\phantomsection\label{fac:III:4:no:69:map:alpha}
For each $n\in\mathbf{Z}$, we defined in
\hyperref[fac:III:3:no:62:def:alpha]{no.~62} a linear map
$\alpha:M_n\longrightarrow H^0(M(n))$; it follows at once that the sum
of these maps defines a homomorphism, again denoted by $\alpha$, from
$M$ to $B^0(M)$.

\phantomsection\label{fac:III:4:no:69:prop:1}
\textsc{Proposition 1.} \textit{Let $k$ be an integer $\geq0$, and let
$J_k$ be the ideal $(t_0^k,\ldots,t_r^k)$ of $S$. For every graded
$S$-module $M$, the graded $S$-modules $B_k^q(M)$ and
$\operatorname{Ext}_S^q(J_k,M)$ are isomorphic.}

\phantomsection\label{fac:III:4:no:69:def:Lk}
Let $L_k^q$, $q=0,\ldots,r$, be the free graded $S$-module with basis
elements $e\langle i_0\cdots i_q\rangle$,
$0\leq i_0<i_1<\cdots<i_q\leq r$, of degree $k(q+1)$. Define an
operator $d:L_k^{q+1}\longrightarrow L_k^q$ and an operator
$\varepsilon:L_k^0\longrightarrow J_k$ by the formulas
\phantomsection\label{fac:III:4:no:69:formula:d}
\[
d(e\langle i_0\cdots i_{q+1}\rangle)
=\sum_{j=0}^{j=q+1}(-1)^j t_{i_j}^k\cdot
e\langle i_0\cdots\widehat{\imath}_j\cdots i_{q+1}\rangle
\]
and
\phantomsection\label{fac:III:4:no:69:formula:augmentation}
\[
\varepsilon(e\langle i\rangle)=t_i^k.
\]

\phantomsection\label{fac:III:4:no:69:lemma:1}
\textsc{Lemma 1.} \textit{The sequence of homomorphisms}
\phantomsection\label{fac:III:4:no:69:formula:resolution-Jk}
\[
0\longrightarrow L_k^r\mathop{\longrightarrow}^{d}L_k^{r-1}
\longrightarrow\cdots\longrightarrow L_k^0
\mathop{\longrightarrow}^{\varepsilon}J_k\longrightarrow0
\]
\textit{is exact.}

For $k=1$ this result is well known (cf.~[6], Chapter~VIII, \S4); the
general case is proved in the same way (or reduced to it). One may also
use the theorem proved in [11].

\hyperref[fac:III:4:no:69:prop:1]{Proposition~1} follows immediately from
\hyperref[fac:III:4:no:69:lemma:1]{the lemma}, on observing that the complex
formed by the
$\operatorname{Hom}_S(L_k^q,M)$ and the transpose of $d$ is precisely the
complex $\sum_{n\in\mathbf{Z}}C_k(M(n))$.

\phantomsection\label{fac:III:4:no:69:cor:1}
\textsc{Corollary 1.} $H_k^q(M)$ \textit{is isomorphic to}
$\operatorname{Ext}_S^q(J_k,M)_0$.

Indeed, these two groups are the degree-zero components of the graded
groups $B_k^q(M)$ and $\operatorname{Ext}_S^q(J_k,M)$.

\phantomsection\label{fac:III:4:no:69:cor:2}
\textsc{Corollary 2.} $H^q(M)$ \textit{is isomorphic to}
$\varinjlim_{k\to\infty}\operatorname{Ext}_S^q(J_k,M)_0$.

It is readily seen that the homomorphism
\hyperref[fac:III:3:no:61:formula:rho]
{$\rho_k^h:H_k^q(M)\longrightarrow H_h^q(M)$} of
\hyperref[fac:III:3:no:61]{no.~61} is carried by
the isomorphism of \hyperref[fac:III:4:no:69:cor:1]{Corollary~1} to the
homomorphism
\phantomsection\label{fac:III:4:no:69:formula:Ext-transition}
\[
\operatorname{Ext}_S^q(J_k,M)_0\longrightarrow
\operatorname{Ext}_S^q(J_h,M)_0
\]
defined by the inclusion $J_h\longrightarrow J_k$; this proves
\hyperref[fac:III:4:no:69:cor:2]{Corollary~2}.

\phantomsection\label{fac:III:4:no:69:remark}
\textsc{Remark.} Let $M$ be a graded $S$-module of finite type. By
\hyperref[fac:II:4:no:48]{no.~48}, $M$ defines a coherent algebraic sheaf
$\mathfrak F'$ on $K^{r+1}$, hence on $Y=K^{r+1}-\{0\}$, and one may
verify that $H^q(Y,\mathfrak F')$ is isomorphic to $B^q(M)$.
